\begin{table}[htbp]
\centering
\caption{Summary Statistics}
\label{tab:summary}
\begin{tabular}{lcccc}
\toprule
 & \multicolumn{2}{c}{Pre-Durbin} & \multicolumn{2}{c}{Post-Durbin} \\
 & \multicolumn{2}{c}{(2005--2010)} & \multicolumn{2}{c}{(2012--2019)} \\
\cmidrule(lr){2-3} \cmidrule(lr){4-5}
 & Mean & (SD) & Mean & (SD) \\
\midrule
Bank Branches & 46.93 & (96.47) & 44.86 & (94.79) \\
Banking Employment & 657.76 & (2051.42) & 656.59 & (2297.22) \\
Average Weekly Wage (\$) & 743.82 & (199.70) & 948.38 & (295.46) \\
Population & 151143.84 & (396004.53) & 158917.70 & (416108.10) \\
Durbin Exposure & 0.35 & (0.27) & 0.35 & (0.27) \\
Branches per 100K & 45.88 & (26.23) & 43.03 & (26.06) \\
Banking Emp. per 100K & 478.47 & (294.66) & 426.28 & (286.48) \\
\midrule
County-years & \multicolumn{2}{c}{10,950} & \multicolumn{2}{c}{12,625} \\
Counties & \multicolumn{4}{c}{2,499} \\
\bottomrule
\end{tabular}
\begin{minipage}{0.95\textwidth}
\vspace{0.5em}
\footnotesize \textit{Notes:} Banking employment and wages from BLS QCEW (NAICS 522110). Branch counts from FDIC Summary of Deposits. Durbin Exposure = share of 2010 county deposits in banks with assets $>$ \$10 billion. Summary statistics exclude the transition year 2011; the full regression sample (25,426 obs) includes 2011. Year 2016 is missing from all panels due to a data download failure. Healthcare sector has fewer observations (18,528) due to QCEW suppression of small county-industry cells.
\end{minipage}
\end{table}

